\begin{figure}[th]
\begin{center}
	\includegraphics[width=0.96\linewidth]{fig/heat.jpg}
\end{center}
\vspace{-8pt}
\caption{\small
\textbf{Token dependency plotted.} The normalized heat map of attention scores in the last self-attention layer of VQGAN encoder is visualized.
4 random 256$\times$256 images from ImageNet validation set are used.
}
\vspace{-2pt}
\label{fig:attn}
\end{figure}

\section{Token dependency in VQVAE} \label{app:dependency}

To examine the token dependency in VQVAE~\cite{vqgan}, we check the attention scores in the self-attention layer before the vector quantization module.
We randomly sample 4 256$\times$256 images from the ImageNet validation set for this analysis.
Note the self-attention layer in \cite{vqgan} only has 1 head so for each image we just plot one attention map.
The heat map in \figref{fig:attn} shows the attention scores of each token to all other tokens, which indicate a strong, bidirectional dependency among all tokens.
This is not surprising since the VQVAE model, trained to reconstruct images, leverages self-attention layers without any attention mask.
Some work~\cite{phenaki} has used causal attention in self-attention layers of a video VAE, but we did not find any image VAE work uses causal self-attention.
% , which could be an interesting direction for future research.
%  is trained to reconstruct the holistic image from all quantized tokens, and 


\section{Time complexity of AR and VAR generation} \label{app:complexity}

We prove the time complexity of AR and VAR generation.
\begin{lemma}
For a standard self-attention transformer, the time complexity of AR generation is $\mathcal{O}(n^6)$, where $h=w=n$ and $h, w$ are the height and width of the VQ code map, respectively.
% A process of AR generation needs $\mathcal{O}(n^2)$, where $n$ is the number of tokens.
\end{lemma}
\begin{proof}
The total number of tokens is $h\times w=n^2$.
For the $i$-th ($1 \le i \le n^2$) autoregressive iteration, the attention scores between each token and all other tokens need to be computed, which requires $\mathcal{O}(i^2)$ time.
% A standard autoregressive generation process needs to compute the attention scores between each token and all other tokens, which requires $\mathcal{O}(n^2)$ time.
So the total time complexity would be:
\begin{align}
	\sum\limits_{i=1}^{n^2} i^2 = \frac{1}{6}n^2(n^2+1)(2n^2+1),
\end{align}
Which is equivalent to $\mathcal{O}(n^6)$ basic computation.
\end{proof}

\noindent For VAR, it needs us to define the resolution sequense $(h_1, w_1, h_2, w_2, \dots, h_K, w_K)$ for autoregressive generation, where $h_i, w_i$ are the height and width of the VQ code map at the $i$-th autoregressive step, and $h_K=h, w_K=w$ reaches the final resolution.
Suppose $n_k = h_k = w_k$ for all $1\le k\le K$ and $n=h=w$, for simplicity.
We set the resolutions as $n_k=a^{(k-1)}$ where $a>1$ is a constant such that $a^{(K-1)}=n$.
\begin{lemma}
For a standard self-attention transformer and given hyperparameter $a>1$, the time complexity of VAR generation is $\mathcal{O}(n^4)$, where $h=w=n$ and $h, w$ are the height and width of the last (largest) VQ code map, respectively.
\end{lemma}
\begin{proof}
Consider the $k$-th ($1 \le k \le K$) autoregressive generation.
The total number of tokens of current all token maps $(r_1, r_2, \dots, r_k)$ is:
\begin{align}
\sum\limits_{i=1}^{k} n_i^2 = \sum\limits_{i=1}^{k} a^{2\cdot(k-1)} = \frac{a^{2k}-1}{a^2-1}.
\end{align}
So the time complexity of the $k$-th autoregressive generation would be:
\begin{align}
	\left(\frac{a^{2k}-1}{a^2-1}\right) ^ 2.
\end{align}
By summing up all autoregressive generations, we have:
\begin{align}
	& \sum\limits_{k=1}^{\log_a(n)+1} \left(\frac{a^{2k}-1}{a^2-1}\right) ^ 2 \\
	&= \frac{(a^4-1)\log n + \left( a^8n^4 - 2a^6n^2 - 2a^4(n^2-1) + 2 a^2 - 1 \right)\log a}{(a^2-1)^3(a^2+1)\log a} \\
	&\sim \mathcal{O}(n^4).
\end{align}
This completes the proof.
\end{proof}



\begin{figure}[th]
\begin{center}
  \includegraphics[width=1\linewidth]{fig/33.jpg}
\end{center}
\vspace{-8pt}
\caption{\small
\textbf{Model comparison on ImageNet 256$\times$256 benchmark.}
More generated 512$\times$512 samples by VAR can be found in the submitted Supplementary Material zip file.
}
\vspace{-2pt}
\end{figure}


\begin{figure}[th]
	\begin{center}
	\includegraphics[width=1\linewidth]{fig/256.jpg}
\end{center}
\vspace{-8pt}
\caption{\small
\textbf{Some generated 256$\times$256 samples by VAR trained on ImageNet.}
More generated 512$\times$512 samples by VAR can be found in the submitted Supplementary Material zip file.
}
\vspace{-2pt}
\end{figure}
